% Options for packages loaded elsewhere
\PassOptionsToPackage{unicode}{hyperref}
\PassOptionsToPackage{hyphens}{url}
\PassOptionsToPackage{dvipsnames,svgnames,x11names}{xcolor}
\documentclass[
  12pt,
  a4paper,
]{report}
\usepackage{xcolor}
\usepackage[margin=2.5cm]{geometry}
\usepackage{amsmath,amssymb}
\setcounter{secnumdepth}{2}
\usepackage{iftex}
\ifPDFTeX
  \usepackage[T1]{fontenc}
  \usepackage[utf8]{inputenc}
  \usepackage{textcomp} % provide euro and other symbols
\else % if luatex or xetex
  \usepackage{unicode-math} % this also loads fontspec
  \defaultfontfeatures{Scale=MatchLowercase}
  \defaultfontfeatures[\rmfamily]{Ligatures=TeX,Scale=1}
\fi
\usepackage{lmodern}
\ifPDFTeX\else
  % xetex/luatex font selection
\fi
% Use upquote if available, for straight quotes in verbatim environments
\IfFileExists{upquote.sty}{\usepackage{upquote}}{}
\IfFileExists{microtype.sty}{% use microtype if available
  \usepackage[]{microtype}
  \UseMicrotypeSet[protrusion]{basicmath} % disable protrusion for tt fonts
}{}
\usepackage{setspace}
\makeatletter
\@ifundefined{KOMAClassName}{% if non-KOMA class
  \IfFileExists{parskip.sty}{%
    \usepackage{parskip}
  }{% else
    \setlength{\parindent}{0pt}
    \setlength{\parskip}{6pt plus 2pt minus 1pt}}
}{% if KOMA class
  \KOMAoptions{parskip=half}}
\makeatother
\usepackage{longtable,booktabs,array}
\usepackage{caption}
\captionsetup[table]{skip=6pt}
\newcounter{none} % for unnumbered tables
\usepackage{calc} % for calculating minipage widths
% Correct order of tables after \paragraph or \subparagraph
\usepackage{etoolbox}
\makeatletter
\patchcmd\longtable{\par}{\if@noskipsec\mbox{}\fi\par}{}{}
\makeatother
% Allow footnotes in longtable head/foot
\IfFileExists{footnotehyper.sty}{\usepackage{footnotehyper}}{\usepackage{footnote}}
\makesavenoteenv{longtable}
\setlength{\emergencystretch}{3em} % prevent overfull lines
\providecommand{\tightlist}{%
  \setlength{\itemsep}{0pt}\setlength{\parskip}{0pt}}
\usepackage{bookmark}
\IfFileExists{xurl.sty}{\usepackage{xurl}}{} % add URL line breaks if available
\urlstyle{same}
% fallback for those not using the hyperref driver hyperxmp:
\makeatletter
\@ifundefined{xmpquote}{\newcommand{\xmpquote}[1]{#1}}{}
\makeatother
\hypersetup{
  pdftitle={COMP2043.GRP Interim Group Report},
  colorlinks=true,
  linkcolor={blue},
  filecolor={Maroon},
  citecolor={blue},
  urlcolor={blue},
  pdfcreator={LaTeX via pandoc}}
% PDF bookmarks (the reader's sidebar outline): open at chapter level when the PDF opens.
\hypersetup{bookmarksopen=true, bookmarksopenlevel=0}
% Figure placement. The default H keeps each figure exactly where it is in the
% Markdown; set figure-placement to e.g. htbp to let LaTeX float figures instead.
\usepackage{float}
\floatplacement{figure}{H}

\ifXeTeX
\special{pdf:trailerid [ <588f9b105b259941974f4ac8475fc568> <588f9b105b259941974f4ac8475fc568> ]}
\fi
\ifPDFTeX
\pdftrailerid{}
\pdftrailer{/ID [ <588f9b105b259941974f4ac8475fc568> <588f9b105b259941974f4ac8475fc568> ]}
\fi
\ifLuaTeX
\pdfvariable trailerid {[ <588f9b105b259941974f4ac8475fc568> <588f9b105b259941974f4ac8475fc568> ]}
\fi

\title{COMP2043.GRP Interim Group Report}
\author{}
\date{}
% Running heads such as "Chapter 3. Requirements", as in the original report.
\pagestyle{headings}

\begin{document}
\pagenumbering{roman}
\begin{titlepage}
\centering
\IfFileExists{figures/team-logo.png}{%
  \includegraphics[width=0.6\linewidth]{figures/team-logo.png}%
}{%
  \fbox{\parbox[c][3cm][c]{0.6\linewidth}{\centering\small Logo placeholder: \texttt{\detokenize{figures/team-logo.png}}}}%
}\par
\vspace{1.5cm}
{\LARGE\bfseries COMP2043.GRP Interim Group Report\par}
\vspace{0.5cm}
Submitted: \textbf{DD/MM/YYYY}\par
\vspace{1.5cm}
{\Large\bfseries Project Title\par}
\vspace{0.3cm}
Supervised by \textbf{Supervisor Name}\par
\vspace{1.5cm}
{\large\bfseries Team N\par}
\vspace{0.3cm}
\textbf{Member Name - 00000000 - username}\par
\textbf{Member Name - 00000000 - username}\par
\vfill
School of Computer Science\par
University of Nottingham\par
\end{titlepage}
\setcounter{page}{2}

{
\hypersetup{linkcolor=blue}
\setcounter{tocdepth}{2}
\tableofcontents
}
\setstretch{1.25}
\cleardoublepage
\pagenumbering{arabic}
\chapter{Introduction}\label{introduction}

Opening paragraph: one-paragraph overview of the problem the project
tackles and what the team will deliver.

\textbf{Keywords:} keyword one, keyword two, keyword three

\section{Current Situation}\label{current-situation}

Describe the project's starting point, the stakeholder situation, and
how the team settled on the project direction.

\section{Motivations and Objectives}\label{motivations-and-objectives}

Explain why the problem matters, then state the project's main
objectives.

\section{Logo Design}\label{logo-design}

Describe the team logo and what its visual elements represent.

\begin{figure}
\centering
\fbox{\parbox[c][4cm][c]{0.8\linewidth}{\centering Placeholder: team logo\\[0.5em]\footnotesize\texttt{\detokenize{figures/logo.png}}}}
\caption{Placeholder: team logo}
\end{figure}

\section{Structure of the Report}\label{structure-of-the-report}

Briefly introduce how the report is organised, then summarise each
chapter:

\begin{itemize}
\tightlist
\item
  \textbf{Introduction:} One-line summary of this chapter.
\item
  \textbf{Background:} One-line summary of this chapter.
\item
  \textbf{Requirements:} One-line summary of this chapter.
\item
  \textbf{Implementation:} One-line summary of this chapter.
\item
  \textbf{User Interface and Prototype:} One-line summary of this
  chapter.
\item
  \textbf{Time Plan:} One-line summary of this chapter.
\item
  \textbf{Summary and Evaluation:} One-line summary of this chapter.
\end{itemize}

\section{Abbreviations}\label{abbreviations}

The following abbreviations are used in this report frequently:

{\def\LTcaptype{none} % do not increment counter
\begin{longtable}[]{@{}ll@{}}
\toprule\noalign{}
Abbreviation & Meaning \\
\midrule\noalign{}
\endhead
\bottomrule\noalign{}
\endlastfoot
ABC & Full term for ABC \\
XYZ & Full term for XYZ \\
\end{longtable}
}

\chapter{Background}\label{background}

\section{Background Introduction}\label{background-introduction}

Outline what this chapter covers.

\section{Development of Driverless
Cars}\label{development-of-driverless-cars}

Summarise the industry's progress and key players, with citations
\hyperref[ref-1]{{[}1{]}}.

\section{Level of Autonomous Driving}\label{level-of-autonomous-driving}

Explain the SAE levels of driving automation and give examples for each.

\section{Testing for Driverless Cars}\label{testing-for-driverless-cars}

\subsection{Introduction}\label{introduction-1}

Introduce the main categories of testing for driverless cars.

\subsection{Software Testing}\label{software-testing}

Describe software testing approaches used in self-driving systems.

\subsection{Simulation Testing}\label{simulation-testing}

Describe simulation-based testing and its benefits.

\subsection{Navigation Layer Testing}\label{navigation-layer-testing}

Describe what the navigation layer does and why it needs testing.

\section{Metamorphic Testing For Driverless
Cars}\label{metamorphic-testing-for-driverless-cars}

\subsection{Oracle Problem}\label{oracle-problem}

Define the test oracle and explain the oracle problem.

\subsection{Concept of Metamorphic
Testing}\label{concept-of-metamorphic-testing}

Explain metamorphic testing and metamorphic relations, with a short
example.

\subsection{Test Case for Autonomous
Driving}\label{test-case-for-autonomous-driving}

Summarise existing work on generating test cases for autonomous driving.

\section{Existing Systems that Address Similar
Problems}\label{existing-systems-that-address-similar-problems}

Summarise real-world incidents and existing systems related to the
problem.

\section{Summary of the Chapter}\label{summary-of-the-chapter}

Recap this chapter and introduce the next one.

\chapter{Requirements}\label{requirements}

\section{Purpose of Project}\label{purpose-of-project}

State what the system is for and who its target users are.

\section{Requirements Elicitation}\label{requirements-elicitation}

Describe the elicitation approach (e.g.~user-centred design).

\subsection{Personas}\label{personas}

Present the user personas, then a short user story for each.

\begin{longtable}[]{@{}lllll@{}}
\caption{Personas}\tabularnewline
\toprule\noalign{}
info / type & persona A & persona B & persona C & persona D \\
\midrule\noalign{}
\endfirsthead
\toprule\noalign{}
info / type & persona A & persona B & persona C & persona D \\
\midrule\noalign{}
\endhead
\bottomrule\noalign{}
\endlastfoot
personal information & \ldots{} & \ldots{} & \ldots{} & \ldots{} \\
occupation & \ldots{} & \ldots{} & \ldots{} & \ldots{} \\
motivation & \ldots{} & \ldots{} & \ldots{} & \ldots{} \\
\end{longtable}

\section{Requirements}\label{requirements-1}

Explain the split between functional and non-functional requirements.

\subsection{Functional Requirements}\label{functional-requirements}

List the requirements for each subsystem.

\begin{itemize}
\tightlist
\item
  User need.

  \begin{itemize}
  \tightlist
  \item
    The system shall \ldots{}
  \item
    The system shall \ldots{}
  \end{itemize}
\end{itemize}

\subsection{Non-Functional
Requirements}\label{non-functional-requirements}

List the quality requirements.

\begin{itemize}
\tightlist
\item
  Quality goal.

  \begin{itemize}
  \tightlist
  \item
    The system should \ldots{}
  \end{itemize}
\end{itemize}

\subsection{UML Diagram}\label{uml-diagram}

Describe what the UML diagram shows.

\begin{figure}
\centering
\fbox{\parbox[c][4cm][c]{0.8\linewidth}{\centering Placeholder: UML use-case diagram of the system\\[0.5em]\footnotesize\texttt{\detokenize{figures/uml-diagram.png}}}}
\caption{Placeholder: UML use-case diagram of the system}
\end{figure}

\section{Requirements Validation}\label{requirements-validation}

Describe how the requirements were validated.

\section{Summary of the Chapter}\label{summary-of-the-chapter-1}

Recap this chapter and introduce the next one.

\chapter{Implementation}\label{implementation}

Introduce the implementation decisions covered in this chapter.

\section{Main Developing Tools for MT to
Apollo}\label{main-developing-tools-for-mt-to-apollo}

Describe the testing platform and build tools, and why the team chose
them.

\begin{figure}
\centering
\fbox{\parbox[c][4cm][c]{0.8\linewidth}{\centering Placeholder: Apollo architecture diagram\\[0.5em]\footnotesize\texttt{\detokenize{figures/apollo-architecture.png}}}}
\caption{Placeholder: Apollo architecture diagram}
\end{figure}

\section{Environment configuration of
Apollo}\label{environment-configuration-of-apollo}

List the prerequisites for installing and running Apollo.

\subsection{Linux System}\label{linux-system}

Which OS version was chosen, and why.

\subsection{NVIDIA GPU DRIVER}\label{nvidia-gpu-driver}

Why the GPU driver is needed and how it was set up.

\subsection{Docker}\label{docker}

What Docker is and which version Apollo requires.

\subsection{The NVIDIA Container
Toolkit}\label{the-nvidia-container-toolkit}

What the toolkit provides and why it was installed.

\subsection{Download and Compile Apollo Source
Code}\label{download-and-compile-apollo-source-code}

Summarise the steps for cloning, building and verifying Apollo.

\begin{figure}
\centering
\fbox{\parbox[c][4cm][c]{0.8\linewidth}{\centering Placeholder: terminal screenshot of a successful Apollo installation\\[0.5em]\footnotesize\texttt{\detokenize{figures/apollo-successful-installation.png}}}}
\caption{Placeholder: terminal screenshot of a successful Apollo
installation}
\end{figure}

\section{OER website Developing
Tools}\label{oer-website-developing-tools}

Describe the languages and database used to build the website.

\section{Other Supporting Tools}\label{other-supporting-tools}

Describe the prototyping, documentation, version control and diagram
tools the team uses.

\section{Software Engineering
Methodologies}\label{software-engineering-methodologies}

Describe the methodologies the team follows (e.g.~Agile, pair
programming).

\section{Summary of the Chapter}\label{summary-of-the-chapter-2}

Recap this chapter and introduce the next one.

\chapter{User Interface and
Prototype}\label{user-interface-and-prototype}

Introduce the prototyping approach and what this chapter shows.

\section{User Interface}\label{user-interface}

List the parts of the user interface covered below.

\subsection{Front Page}\label{front-page}

Describe the front page layout and its interactive elements.

\begin{figure}
\centering
\fbox{\parbox[c][4cm][c]{0.8\linewidth}{\centering Placeholder: front page mock-up\\[0.5em]\footnotesize\texttt{\detokenize{figures/front-page.png}}}}
\caption{Placeholder: front page mock-up}
\end{figure}

\subsection{Login and Registration
Page}\label{login-and-registration-page}

Describe the login and registration flow.

\begin{figure}
\centering
\fbox{\parbox[c][4cm][c]{0.8\linewidth}{\centering Placeholder: account page mock-up\\[0.5em]\footnotesize\texttt{\detokenize{figures/account-page.png}}}}
\caption{Placeholder: account page mock-up}
\end{figure}

\subsection{Search Box Page}\label{search-box-page}

Describe the search function and its result pages.

\begin{figure}
\centering
\fbox{\parbox[c][4cm][c]{0.8\linewidth}{\centering Placeholder: search result page mock-up\\[0.5em]\footnotesize\texttt{\detokenize{figures/search-result-page.png}}}}
\caption{Placeholder: search result page mock-up}
\end{figure}

\section{Prototype}\label{prototype}

List the prototype pages covered below.

\subsection{Background Page}\label{background-page}

Describe the background page.

\begin{figure}
\centering
\fbox{\parbox[c][4cm][c]{0.8\linewidth}{\centering Placeholder: background page mock-up\\[0.5em]\footnotesize\texttt{\detokenize{figures/background-page.png}}}}
\caption{Placeholder: background page mock-up}
\end{figure}

\subsection{Metamorphic Testing Page}\label{metamorphic-testing-page}

Describe the testing pages, the interactive test visualisation and the
quiz.

\begin{figure}
\centering
\fbox{\parbox[c][4cm][c]{0.8\linewidth}{\centering Placeholder: testing page mock-up\\[0.5em]\footnotesize\texttt{\detokenize{figures/testing-page.png}}}}
\caption{Placeholder: testing page mock-up}
\end{figure}

\begin{figure}
\centering
\fbox{\parbox[c][4cm][c]{0.8\linewidth}{\centering Placeholder: quiz page mock-up\\[0.5em]\footnotesize\texttt{\detokenize{figures/quiz-page.png}}}}
\caption{Placeholder: quiz page mock-up}
\end{figure}

\subsection{Forum Page}\label{forum-page}

Describe the forum's posting, commenting and liking features.

\begin{figure}
\centering
\fbox{\parbox[c][4cm][c]{0.8\linewidth}{\centering Placeholder: forum page mock-up\\[0.5em]\footnotesize\texttt{\detokenize{figures/forum-page.png}}}}
\caption{Placeholder: forum page mock-up}
\end{figure}

\subsection{Others Page}\label{others-page}

Describe the team information and contact pages.

\begin{figure}
\centering
\fbox{\parbox[c][4cm][c]{0.8\linewidth}{\centering Placeholder: team information page mock-up\\[0.5em]\footnotesize\texttt{\detokenize{figures/team-information-page.png}}}}
\caption{Placeholder: team information page mock-up}
\end{figure}

\begin{figure}
\centering
\fbox{\parbox[c][4cm][c]{0.8\linewidth}{\centering Placeholder: contact us page mock-up\\[0.5em]\footnotesize\texttt{\detokenize{figures/contact-us-page.png}}}}
\caption{Placeholder: contact us page mock-up}
\end{figure}

\section{Summary of the Chapter}\label{summary-of-the-chapter-3}

Recap this chapter and outline the next steps for the prototype.

\chapter{Time Plan}\label{time-plan}

\section{Schedule}\label{schedule}

Present the Gantt chart, then describe each work stream (preparation and
design, implementation phases) and who is responsible.

\begin{figure}
\centering
\fbox{\parbox[c][4cm][c]{0.8\linewidth}{\centering Placeholder: project Gantt chart\\[0.5em]\footnotesize\texttt{\detokenize{figures/project-gantt-chart.png}}}}
\caption{Placeholder: project Gantt chart}
\end{figure}

\section{Deliverables}\label{deliverables}

List the expected outputs:

\begin{itemize}
\tightlist
\item
  Deliverable one
\item
  Deliverable two
\end{itemize}

\section{Summary of the Chapter}\label{summary-of-the-chapter-4}

Recap the timeline and the main deliverables.

\chapter{Summary and Evaluation}\label{summary-and-evaluation}

\section{Problems Encountered and
Reflections}\label{problems-encountered-and-reflections}

\subsection{Problems in Design}\label{problems-in-design}

Describe the design and technical problems, how the team handled them,
and what was learned.

\begin{figure}
\centering
\fbox{\parbox[c][4cm][c]{0.8\linewidth}{\centering Placeholder: user-centred design process diagram\\[0.5em]\footnotesize\texttt{\detokenize{figures/ucd-process.png}}}}
\caption{Placeholder: user-centred design process diagram}
\end{figure}

\subsection{Team Management Issues}\label{team-management-issues}

Describe the team management problems and how the team overcame them.

\section{Conclusion}\label{conclusion}

Summarise what the team achieved this semester and the plan going
forward.

\chapter*{Bibliography}\label{bibliography}
\addcontentsline{toc}{chapter}{Bibliography}

\begin{enumerate}
\def\labelenumi{\arabic{enumi}.}
\tightlist
\item
  \protect\phantomsection\label{ref-1}{}Author, A. \emph{Title of the
  web page.} \url{https://example.com}, Year. {[}Online; accessed
  DD-Month-YYYY{]}.
\item
  \protect\phantomsection\label{ref-2}{}Author, A., and Author, B. Title
  of the article. \emph{Journal Name 1}, 2 (Year), 1--10.
  doi:10.0000/example.
\end{enumerate}

\appendix

\chapter{Meeting Minutes}\label{meeting-minutes}

\section{Meeting 1}\label{meeting-1}

\begin{itemize}
\tightlist
\item
  \textbf{Chairperson:} Name
\item
  \textbf{Recorder:} Name
\item
  \textbf{Date:} Month DD, YYYY, 9:00 a.m. - 9:50 a.m.
\item
  \textbf{Place:} Room or Online
\end{itemize}

Summarise the topics discussed, guidance for next week, and the next
meeting's info.

\section{Meeting 2}\label{meeting-2}

Same structure as Meeting 1.

\end{document}
